\paragraph{Reduction of the marked Wach lattices.}
\label{par:residue-wach-realization}
Use the field, coefficient extension and reduced cyclotomic
realization of Section~\ref{subsec:cut-wach-models} with
\(p=\ell\), \(u=1+p\), and \(X\) the reduction of the integral
parameter \(X_0=-Z\). The residue field is \(k'\).
All crystalline and residual factors use the terminal Hodge
labels of that section.
The residual parameter actions, with Frobenius from the preceding letter, are
\[
 \varphi(X)=X^p,\qquad \gamma(X)\equiv X-X^2\pmod{X^3}.
\]


The first scalar character \(\chi_1\) has jump \(w_i\);
the scalar characters are \(\chi_c=\chi_1\epsilon^{1-c}\)
for \(1\le c\le t\), with \(\epsilon\) the cyclotomic character.
A full representation is crystalline with the labeled jumps
of the marked rank block together with those of the scalar
characters. Its residual characteristic polynomials agree with
those of that direct sum on the Galois group.

For a stable lattice
in such a full sum, take the Wach lattice after twisting by
\(\chi_1^{-1}\). With \(U_i\) the multiset of rank-block jumps before the twist, the full twisted multiset is
\[
 B_i=(U_i-w_i)\sqcup\{0,\ldots,t-1\}.
\]
Both parts retain their multiplicities. The separate rank and line lattices
are treated by the same twist. Their labeled height bounds,
after reduction by the coefficient uniformizer, give
\[
      X^{w_i^*}N_i\ \subset\ \Phi_i N_{i-1}
                    \ \subset\ X^{-w_i}N_i,
              \qquad (\gamma_i-1)N_i\subset XN_i .
\]
Here \(N_i\) is the reduction of the tame-invariant Wach lattice
inside the residual 'etale module on the corresponding letter.
A scalar line of sole jump \(e\) has Frobenius ideal \(X^e\).
For the trivial quotient line the invariant basis is \(1\).
\cite{berger-wach}\label{cor:residue-wach-realization-1}

\paragraph{The polynomial on a high-end letter.}
\begin{lemma}[Simple zero factor on a high-end letter]
\label{lem:residue-high-end-polynomial}
For the upper-end letter and \(s=w_i^*+1\), write \(D\) for the divided action on \(N_i/X^sN_i\). For \(P\) as in the proof the minimal polynomial satisfies
\[
 m_D(T)\mid TP(T),\qquad P(0)\ne0,\qquad \ker(D^j)=\ker(D)\quad(j\ge1).
\] Modulo \(X\)
the characteristic polynomial still records every jump residue
with its multiplicity.
\end{lemma}

\begin{proof}
In the integral tame parameter \(X_0\), with reduction \(X\),
put \(q=X_0+p\) and \(\nu=(\gamma-1)/X_0\). This is
\(O'\)-linear and preserves the invariant Wach lattice, since
\(X_0\) generates the invariant augmentation ideal and the
action is trivial on constants. For \(s\ge1\) define
\[
 I_s=\{a-b:\ 0\le a<s,\ b\in B_i\},\qquad
 d(m)=\frac{1-(1+p)^m}{p}\in\mathbb Z_p .
\]
On the characteristic-zero quotient by \(q^s\), the operator
is killed by
\[
 \prod_{m\in I_s}(T-d(m)),
\]
using each distinct
integer once. At the primitive level choose crystalline directions
adapted to the filtration. In a direction of jump \(b\) the
quotient has Laurent degrees \(-b,\ldots,s-1-b\) in the
local logarithmic parameter. The scalar substitution multiplies
this parameter by \(1+p\), and \(X_0=-p\) at the root.

In a direction of jump \(b\) the triangular operator \(T\) has distinct diagonal entries \(d(a-b)\), hence
\[
 \prod_{a=0}^{s-1}\big(T-d(a-b)\big)=0
 \quad\text{on this direction}.
\]
The product over the union of the eigenvalues thus kills every direction, including repeated jumps.
This is the finite-level filtration comparison of
Paragraph~\ref{par:wach-level-models}. The quotient of the
free invariant Wach lattice by \(q^s\) is free over \(O'\)
by monic division. The polynomial identity therefore holds
integrally. For \(s=1\) the characteristic polynomial has each
\(d(-b)\) with its multiplicity in \(B_i\), by the same
calculation and injectivity on coefficients.
\label{cor:residue-wach-realization-2}

Reduction gives \(q=X\), \(\bar d(m)=-m\), and operator
\[
 \nabla=(\gamma_i-1)/X.
\]
For negative \(m\) the congruence
for \(d(m)\) follows by inverting \((1+p)^{-m}\) modulo \(p^2\)
before division. The characteristic polynomial modulo \(X\)
therefore counts the residues of \(B_i\) with multiplicity.
At the upper end take \(s=w_i^*+1\). If \(0\le a<s\) and
\(a\equiv b\bmod p\), then \(a=b\). For a scalar or central
jump, writing \(b=u-w_i\), both \(u\) and \(w_i+a\) are in
the interval of width \(\eta'_i\) next to \(h_i\), including
the endpoint \(h_i\) for the latter.

For an outer jump the
congruence would give \(h_i-u\equiv w_i^*-a\bmod p\) in the
forbidden end window. It follows that
\[
   P(T)=\prod_{m\in I_s\setminus\{0\}}(T+m)\in k'[T],
   \qquad P(0)\ne0,\qquad
   \nabla P(\nabla)N_i\subset X^sN_i .
\]
If zero is absent from \(I_s\), \(P(\nabla)\) already has
image in \(X^sN_i\); applying \(\nabla\) gives the
displayed identity since the divided action preserves that
lattice. In the finite quotient, \(P(\nabla)\) is invertible
on the generalized zero space, since \(P(0)\ne0\). Thus
\(\nabla P(\nabla)=0\) there forces \(\nabla=0\). Using
the characteristic polynomial with all multiplicities would
lose this conclusion. \label{cor:residue-wach-realization-3}
\end{proof}

Finally an admissible residual representation is a quotient
of the reduction of its presenting full lattice. Its presenting
map is onto and kills \(\lambda T\), and untwisting preserves
that map. More generally take the intersection with a vector
subspace for a subrepresentation and the image for a quotient.
The intersection lattice is saturated. In a basis adapted to it,
Frobenius base change is exact, since the series ring is free
over its Frobenius image on \(1,X,\ldots,X^{p-1}\). Intersecting
the height inclusions or taking images gives both inclusions
for the induced lattices. The divided cyclotomic differences
intertwine the maps, so their polynomials restrict and descend:
for a stable vector subspace \(W\) one has
\[
 (X^sN_i)\cap W=X^s(N_i\cap W).
\]
The induced lattice sequence stays
exact modulo \(X\). The integral Wach intersection gives the
direct sum for a sum of lattices.

For the standard reference
use the split sum of the chosen rank lattice and the character
lines. Its reductions give the required models on those
summands, with their original labels. The prescribed residual
characteristic polynomials fix the composition factors of
each full sum, and untwisting multiplies them all by the
same character. This supplies the models used below with
their two heights, polynomials and target Sen counts.
\label{cor:residue-wach-realization-4}
